\chapter{Síntesis de diseño para la fase 1}
\section{Propuesta de diseño experimental}

Para implementar y comparar los algoritmos, se recomienda seguir los siguientes pasos:

\begin{enumerate}
    \item Definir el HCP como un problema de decisión sobre grafos no dirigidos.
    \item Utilizar permutaciones cíclicas como representación principal.
    \item Validar que cada candidata tenga una estructura de permutación antes de evaluar sus conexiones; permitir conexiones ausentes como penalizaciones durante la búsqueda.
    \item Para permutaciones completas, contar conexiones ausentes, incluido el cierre. Si se usan otras codificaciones, definir también restricciones de cobertura y conectividad.
    \item Implementar SA, ACO, ABC, DE y PSO bajo la misma interfaz de evaluación.
    \item Utilizar el FHCP Challenge Set como benchmark.
    \item Medir tasa de éxito, tiempo, número de evaluaciones, iteraciones y desviación estándar.
    \item Verificar cada ciclo encontrado antes de reportarlo como solución válida.
    \item Comparar los algoritmos bajo el mismo presupuesto computacional.
\end{enumerate}

La fortaleza del diseño está en separar tres capas:

\begin{itemize}
    \item Representación.
    \item Búsqueda.
    \item Evaluación.
\end{itemize}

Esto permite comparar los cinco algoritmos bajo condiciones equivalentes y obtener resultados reproducibles sobre el FHCP Challenge Set. \cite{de2}


\section{Precisiones de la integración}
Se adopta la separación entre representación, búsqueda y evaluación propuesta en investigacion.tex. El decodificador transforma el estado interno en una ruta; el fitness orienta la búsqueda; un verificador independiente decide si la ruta satisface exactamente la definición de ciclo.

OX y PMX son cruces especializados para permutaciones. Intercambio, inserción e inversión conservan los vértices, pero pueden romper conexiones. Si se codifican aristas, exigir grado 2 en cada vértice no basta: pueden aparecer varios ciclos disjuntos; también debe exigirse conectividad.\cite{aco4}

Se corrige el objetivo combinado de violaciones y longitud sugerido en investigacion.tex: toda permutación circular tiene $n$ conexiones candidatas, así que su longitud combinatoria no discrimina calidad. El costo de conexiones ausentes sí distingue incumplimientos. Una lista con duplicados no debe evaluarse como si fuera una permutación válida.

ACO puede construir caminos parciales; su calidad interna no equivale al costo de una permutación completa. La comparación entre algoritmos usa el mismo éxito verificado, y cualquier calidad parcial declara si la ruta es completa o parcial.

\section{Precisión, recall y F1: límites de aplicación}
Para clasificación con referencia independiente,
\[
P=\frac{TP}{TP+FP},\qquad R=\frac{TP}{TP+FN},\qquad
F_1=\frac{2PR}{P+R}=\frac{2TP}{2TP+FP+FN}.
\]
Si un denominador es cero, se informa que la métrica no está definida, salvo que se declare una convención de reporte.\cite{de9}

Se matiza investigacion.tex: agotar el tiempo sin encontrar ciclo es un resultado inconcluso, no una prueba de que el grafo sea negativo. FHCP contiene solo positivos. Si ``detección'' significa encontrar un ciclo verificado, el recall coincide con la tasa de detección, la precisión vale 1 cuando hay detecciones y F1 es una transformación del recall. No son tres pruebas independientes de rendimiento. Evaluar un clasificador de hamiltonicidad exigiría también negativos certificados y una política para resultados inconclusos. Comparar aristas contra un único tour penalizaría otros ciclos correctos.

\section{Registro y reproducibilidad}
Se propone elegir las instancias antes de experimentar y separar grafos de calibración de los de evaluación. Registrar por ejecución: identificador, $n$, $m$, algoritmo, versión, parámetros, semilla, presupuesto, evaluaciones, iteraciones, tiempo, mejor costo y estado final. Guardar cada ciclo como TOUR y verificarlo independientemente.

Deben documentarse equipo y entorno, distinguir tiempo de CPU del tiempo de pared e incluir el costo de decodificación, reparación, búsqueda local y reinicios. Igualar generaciones no iguala esfuerzo: una generación poblacional puede evaluar muchas candidatas. Las evaluaciones deben acompañarse de tiempo, pues las operaciones internas difieren entre métodos.\cite{de5,de6}

Reportar tasa de éxito por instancia y cobertura del conjunto por separado; añadir incertidumbre y dispersión. No descartar fallos al calcular promedios sin explicarlo. El esfuerzo total dividido entre el número de éxitos estima esfuerzo hasta éxito bajo reinicios; con cero éxitos no produce una estimación finita. Las propuestas de benchmarking se transfieren desde optimización y no se presentan como reglas oficiales del FHCP.\cite{de5,de6}

\section{Decisiones pendientes}
Antes de implementar se acuerdan instancias, presupuestos, repeticiones, calibración y uso de búsqueda local. El texto conserva GA como fundamento del módulo evolutivo; los cinco algoritmos solicitados son SA, ACO, ABC, DE y PSO. Todos los ejemplos son didácticos y no prueban superioridad de un método.
